\section{课堂收束与 Q\&A：成本、规模和未解问题}

课堂最后把前述实验压缩成 takeaways，并在 Q\&A 中补充 vendor 成本、pair-order mitigation、reward model scoring 与样本规模的工程直觉。这些信息很有价值，但必须按“讲者现场估计”而不是定理来记录。

\subsection{Takeaways：Recipe 是一个四维对象}

前面各章分别讨论 data、objective 与 evaluator，本节把它们重新压缩成一个可执行的系统模板，并检查讲者最终 bullets 是否覆盖全部环节。读 takeaway 页时不要逐条背诵，而要沿“数据如何进入训练、训练如何被测量、测量如何反馈到下一轮”建立闭环。一份可复现 chatbot recipe 至少应包含 data distribution、optimization stage、evaluation suite 与 provenance audit；若缺任何一维，模型分数就难以解释或复现。

\lecturefigure{slide-66.jpg}{讲者总结：数据、SFT 结果、benchmark gap 与 GPT-4 quirks}{官方幻灯片 p.66}

\begin{knowledgebox}{读图：把 bullets 重排成因果链}
Data amount/length/task/human role 决定训练分布；SFT 与 DPO 改变 policy；TruthfulQA、MT-Bench 等观察不同切面；RM/red-team benchmark 仍有缺口；GPT-4 judge 又受 data/style/position 影响。最后一组 quirks 会反向限制前面所有 leaderboard claim。
\end{knowledgebox}

\subsection{Public Labor 与 Institutional Context}

上一小节完成了技术总结，本节保留讲者刻意放在结尾的社会与团队语境，回答“谁生产数据、谁承担治理、谁提供复现资源”。读这组页面时应把它们视为 provenance 与责任边界，而不是新的 benchmark evidence。课堂用新闻报道、UN advisory body 与 H4 resources 收尾；这些页面不是性能证据，却提醒 alignment recipe 背后有人类劳动、治理与开源协作，省略它们会把系统误写成纯算法流程。

\lecturefigure{slide-67.jpg}{New York Times 对 RLHF human labor 的报道}{官方幻灯片 p.67}

\begin{knowledgebox}{读图：Alignment 的“secret ingredient”包含 labor system}
Data guidelines、vendor management、质量审核和 annotator welfare 都影响最终模型。把它们隐藏为一个 dataset name，会让可复现性、成本和责任边界同时消失。
\end{knowledgebox}

\lecturefigure{slide-68.jpg}{讲者加入 UN AI Advisory Body 的课堂时间背景}{官方幻灯片 p.68}

\begin{knowledgebox}{读图：Institutional role 说明问题超出 benchmark}
Chatbot helpfulness/safety 牵涉全球语言、劳动、治理与公共利益。该页记录 2023 课堂背景，不用于证明任何技术结论；它提醒我们 evaluation target 最终来自社会选择。
\end{knowledgebox}

\lecturefigure{slide-69.jpg}{H4 red teaming、leaderboard 与 dialogue-agent 官方资源}{官方幻灯片 p.69}

\begin{knowledgebox}{读图：三个链接对应三种复现入口}
Red teaming 资源面向漏洞发现，LLM leaderboard 面向 evaluator audit，dialog agents 面向训练实践。读者应从原始方法和数据说明复现实验，而不是只抄榜单数字。
\end{knowledgebox}

\lecturefigure{slide-70.jpg}{H4 核心团队与开源社区致谢}{官方幻灯片 p.70}

\begin{knowledgebox}{读图：研究结果是团队与社区产物}
数据、训练、评估和发布由不同成员共同完成。保留 attribution 不只是礼貌，也帮助追踪代码、数据、paper 与 blog 的来源，避免把所有判断错误归到单一讲者或模型。
\end{knowledgebox}

\subsection{Q\&A：成本与 Scale 只能按单位解释}

Rajani 在 Q\&A 中估计，向 Surge/Scale AI 购买约 20K preference dialogues 与约 10K demonstrations 的总成本约为 50 万美元。这个数字依赖 2023 vendor、任务、轮数和质量流程；它的教学意义是人工数据成本足以改变研究路线，而不是给出当前市场报价。

\teachervoice{讲者比较 synthetic generation 与 vendor data 时说，前者更 cost-effective；human collection 大约花费 half a million dollars。这里的单位分别是 dialogues、prompts 与 demonstrations，不能混成“30K samples”。}

Pair-order bias 的实用 mitigation 是让每个 model 同等概率出现在左右两侧，并在不同 prompts 上生成 reversed ordering。它降低一阶 position effect，却不解决 judge style bias、task verification 和 data coupling。

\teachervoice{面对“哪些 prompts 造成 bias”的提问，讲者坦言团队没有系统研究。她区分了 observed left bias 与 speculative cause，并把 balanced ordering 作为工程缓解，而不是宣称已找到机制解释。}

最后一个问题追问 reward model 如何给 RL policy 信号。课堂用“对完整 sequence 给分、反复采样优化”的直觉解释，并对比约 100K RLHF examples 与约 10K SFT examples 的数量级。

\teachervoice{10K/100K 是讲者现场的 order-of-magnitude intuition：preference/RL pipeline 需要大量 sampled comparisons，SFT demonstration 更直接。它不能脱离 model、task、turn count 与 label unit 写成 sample-complexity theorem。}

\subsection{本章小结}

Recipe 的真实成本包括人类劳动、数据治理与 evaluator validation。Q\&A 提供了规模直觉和未解问题，也再次证明所有数字必须连同单位、协议和证据级别一起记录。

